\documentclass[10pt]{article}

\usepackage[utf8]{inputenc}
\usepackage[T1]{fontenc}
\usepackage{lmodern}
\usepackage{amsmath,amssymb,amsthm}
\usepackage{array,booktabs}
\usepackage[table]{xcolor}
\usepackage{enumitem}
\usepackage{geometry}
\usepackage{hyperref}
\usepackage{algpseudocode}
\usepackage{float}
\usepackage[ruled, vlined, algoruled]{algorithm2e} % Options for styling
\usepackage{tikz}
\geometry{top=0.65in,bottom=0.65in,left=1.2in,right=1.2in}
\hypersetup{
  colorlinks=true,
  linkcolor=blue!60!black,
  urlcolor=cyan!50!black,
  linktoc=all
}
\setlist{nosep,leftmargin=*}
\renewcommand{\arraystretch}{1.2}
\newcolumntype{P}[1]{>{\raggedright\arraybackslash}p{#1}}

\theoremstyle{definition}
\newtheorem{definition}{Definition}[section]
\newtheorem{theorem}[definition]{Theorem}
\newtheorem{lemma}[definition]{Lemma}
\newtheorem{remark}[definition]{Remark}
\newtheorem{corollary}[definition]{Corollary}

\newcommand{\OPT}{\operatorname{OPT}}
\newcommand{\ALG}{\operatorname{ALG}}
\newcommand{\poly}{\operatorname{poly}}
\newcommand{\YES}{\textsc{Yes}}
\newcommand{\NO}{\textsc{No}}
\newcommand{\NP}{\mathsf{NP}}
\newcommand{\coNP}{\mathsf{coNP}}
\newcommand{\PTIME}{\mathsf{P}}
\newcommand{\NPC}{\mathsf{NP\text{-}complete}}

% Colours used to separate the three layers of approximation proofs.
\definecolor{algblue}{RGB}{0,82,155}
\definecolor{optred}{RGB}{180,45,45}
\definecolor{bridgepurple}{RGB}{105,55,145}

\title{CS4234: Optimization Algorithms}
\author{@Wxy2003-xy}
\date{}

\begin{document}
\section{Randomization and Derandomization}\label{ra:main}

\subsection{The indicator-variable proof template}
For a maximization objective with nonnegative local rewards, write
\[
  \ALG=\sum_i w_i I_i,\qquad
  \mathbb E[\ALG]=\sum_i w_i\Pr[I_i=1].
\]
Linearity of expectation does \emph{not} require the indicators to be
independent. Independence is needed only when a probability calculation
multiplies probabilities of separate random choices. If every reward is
collected with probability at least $\alpha$ and
$W=\sum_iw_i\ge\OPT$, then $\mathbb E[\ALG]\ge\alpha\OPT$.
This guarantees the existence of a good outcome, not that every run is good.
An expected approximation bound alone is not a high-probability bound.

\subsection{Weighted MAX-CUT, MAX-$k$-CUT, and MAX-DICUT}
Assume nonnegative edge weights and remove self-loops, which can never cross
a cut. For undirected MAX-$k$-CUT, independently label every vertex uniformly
from $\{1,\ldots,k\}$, for $k\ge2$. Count edges whose endpoint labels differ.
For every edge $e=\{u,v\}$,
\[
  \Pr[I_e=1]=1-\sum_{a=1}^k\Pr[u=a,v=a]=1-\frac1k.
\]
Consequently
\[
  \mathbb E[\ALG]=(1-1/k)W\ge(1-1/k)\OPT.
\]
The ordinary MAX-CUT case is $k=2$, giving $\OPT/2$; MAX-3-CUT gives
$2\OPT/3$. A run takes $O(n+m)$ time including evaluation.

For directed MAX-DICUT, select each vertex into $S$ independently with
probability $1/2$, and count only arcs from $S$ to $V\setminus S$.
Each arc $(u,v)$ is counted with probability $1/4$, so
$\mathbb E[\ALG]=W/4\ge\OPT/4$.
Do not count arcs in both directions: doing so changes the objective.
LP rounding improves the retained fraction to $1/2$
(Section~\ref{lp:dicut}).

\subsection{Conditional expectation: how to remove randomness}\label{ra:conditional}
Let $F$ be the expected final objective conditional on the decisions already
fixed. If the next random variable has outcomes $a$ with probabilities $p_a$,
then
\[
  F=\sum_a p_a F_a,\qquad
  F_a=\mathbb E[\ALG\mid\text{fixed decisions},X=a].
\]
At least one positive-probability branch has $F_a\ge F$. Fix such a branch and
repeat. The invariant is that the conditional expectation never decreases;
after the last decision it equals the actual objective. Thus the deterministic
output is at least the initial expectation. To claim a polynomial-time
algorithm, show how to compute each $F_a$ in polynomial time.

\subsubsection*{MAX-$k$-CUT: computable conditional contributions}
Under a partial labelling, an edge with both endpoints fixed contributes its
full weight if their labels differ and zero otherwise. An edge with at least
one endpoint unfixed contributes $(1-1/k)w_e$. When assigning a vertex $v$,
only edges to already fixed neighbors have label-dependent contributions.
Choose a label $a$ maximizing
\[
  \sum_{u:\{u,v\}\in E,\;u\text{ fixed},\;\operatorname{label}(u)\ne a}w_{uv},
\]
equivalently minimizing the weight of fixed neighbors with label $a$.
Process vertices in any order. This gives a deterministic
$(1-1/k)$ retained-value approximation in $O(m+nk)$ time by accumulating
neighbor weights in $k$ buckets at each vertex.

\textbf{Worked MAX-CUT example.}
For a unit-weight triangle the initial expectation is $3/2$. Fix the first
vertex to side 0; the expectation remains $3/2$. Putting the second on side 1
makes their edge contribute 1 and the two unresolved edges contribute $1/2$
each, giving conditional expectation 2. Either choice for the third vertex
then yields cut value 2.

\subsubsection*{MAX-SAT: computable conditional contributions}
First remove duplicate literals, discard empty clauses (always false), and
separate tautologies (always true). Thus every remaining clause contains
distinct variables and no variable together with its negation. Suppose an
unfixed variable $x_j$ is independently true with probability $p_j$.
For each clause $C_i$, its conditional expected contribution is
\[
  \begin{cases}
    w_i, &\text{a fixed literal is already true},\\
    0, &\text{all literals are fixed and false},\\
    w_i\left(1-\displaystyle\prod_{\ell\text{ unfixed in }C_i}
       q_\ell\right), &\text{otherwise},
  \end{cases}
\]
where $q_{x_j}=1-p_j$ and $q_{\neg x_j}=p_j$ are literal-failure
probabilities. For fair variables and $r$ unfixed literals, the last line is
$w_i(1-2^{-r})$. Evaluate the total for setting the next variable to 0 and
to 1, and keep the larger value. Scanning all $L$ literal occurrences for
each of $n$ variables gives $O(nL)$ arithmetic operations. This works for
fair, biased, and LP-derived independent probabilities.

\textbf{Worked example.}
For unit-weight clauses $(x)$ and $(\neg x\lor y)$ with fair variables,
$F=1/2+3/4=5/4$. Setting $x=1$ gives $F_1=1+1/2=3/2$; setting $x=0$
gives $F_0=0+1=1$. Choose $x=1$, then $y=1$, obtaining value 2.

For a directed cut with arbitrary independent membership probabilities,
an arc's conditional contribution is $w_{uv}p'_u(1-p'_v)$, where $p'_v$ is
0 or 1 for a fixed vertex and its original probability otherwise.
This also makes both MAX-DICUT rounding schemes below derandomizable.

\subsection{Weighted MAX-SAT by fair and biased assignment}\label{ra:sat}
Use the clause preprocessing above; handle tautological clauses as a constant
reward added to every solution. All weights are nonnegative. For a remaining
clause $C_i$ of length $\ell_i\ge1$, fair independent assignment gives
\[
  \Pr[C_i\text{ satisfied}]=1-2^{-\ell_i},\qquad
  \mathbb E[\ALG]=\sum_iw_i(1-2^{-\ell_i})\ge W/2\ge\OPT/2.
\]
The satisfaction probability \emph{increases} with clause length. If every
clause has at least $k$ literals, the guarantee strengthens to
$(1-2^{-k})\OPT$. An assignment and its value can be computed in $O(n+L)$
time, where $L$ is the number of literal occurrences.

\subsubsection*{Biased assignment and the unit-clause issue}
Merge repeated unit clauses by adding their weights. For each variable,
rename it consistently throughout the formula if needed so that the total
positive unit-clause weight is at least the negative one. Let $N$ be the sum
of the negative unit-clause weights after this renaming. Every assignment
loses at least the smaller weight in each opposing pair, so
$\OPT\le W-N$.

Set every renamed variable independently true with probability
$p\in[1/2,1]$. A clause with $a$ positive and $b$ negative literals is
satisfied with probability
\[
  1-(1-p)^a p^b\ge1-p^{a+b}.
\]
In a lower bound, ignore the negative unit clauses. Every remaining unit
clause is satisfied with probability $p$, and every nonunit clause with
probability at least $1-p^2$. Therefore
\[
  \mathbb E[\ALG]\ge\min\{p,1-p^2\}(W-N)
    \ge\min\{p,1-p^2\}\OPT.
\]
The maximum occurs when $p=1-p^2$, giving
\[
  p=\frac{\sqrt5-1}{2},\qquad
  \mathbb E[\ALG]\ge\frac{\sqrt5-1}{2}\OPT.
\]
If no negative unit clauses remain, this simplifies to $N=0$.
Without the renaming/upper-bound argument, assuming all unit clauses are
positive is not justified. Conditional expectation gives deterministic
versions of these guarantees.

\section{Linear Programming and Rounding}\label{lp:main}

\subsection{Modelling, relaxation, and the direction of the bound}\label{lp:template}
An ILP uses integer variables; a $0$--$1$ formulation uses binary decision
variables. A relaxation replaces $x\in\{0,1\}$ by $0\le x\le1$ and keeps
the linear constraints. Its feasible set contains all integral feasible
points. Write $LP^*$ for the \emph{optimal} relaxed objective:
\[
  \begin{array}{c|c|c}
    &\text{relaxation bound}&\text{target rounding proof}\\\hline
    \text{minimization}&LP^*\le\OPT&\ALG\le\rho LP^*\le\rho\OPT\\
    \text{maximization}&LP^*\ge\OPT&
       \mathbb E[\ALG]\ge\alpha LP^*\ge\alpha\OPT
  \end{array}
\]
An arbitrary feasible fractional objective is not necessarily a valid bound
on $\OPT$ in the displayed direction. LP is solvable in polynomial time in
its encoding length; general ILP is NP-hard. After rounding, prove
\emph{feasibility} separately from the objective bound.

\textbf{Proving a model is exact.}
Show both (i) each original solution gives a feasible integral point of the
same objective and (ii) an integral point decodes into a solution of at least
its objective for maximization (at most for minimization). These establish
equality of optimal values. An auxiliary reward variable often needs only
upper bounds: maximizing a nonnegative weighted objective pushes it up when
possible. Such constraints need not force its intended Boolean meaning at
\emph{every feasible point}. If its coefficient is zero, it need not be
pushed up even at an optimum; the decoded solution and optimal value remain
correct. Add equivalence constraints if the variable's exact meaning is
required independently of the objective.

\subsection{Weighted vertex cover: threshold rounding}\label{lp:vc}
For nonnegative vertex costs $c_v$, let $x_v$ indicate selection. The ILP is
\[
  \min\sum_{v\in V}c_vx_v
  \quad\text{subject to}\quad
  x_u+x_v\ge1\quad(\{u,v\}\in E),\qquad x_v\in\{0,1\}.
\]
Relax to $[0,1]$, obtain an optimum $x^*$, and return
$C=\{v:x_v^*\ge1/2\}$. At least one endpoint of every edge reaches $1/2$,
so $C$ is feasible. Also
\[
  c(C)=\sum_{v\in C}c_v
    \le\sum_{v\in C}2c_vx_v^*
    \le2LP^*\le2\OPT.
\]
The algorithm is polynomial time (LP plus a linear scan).
On the unweighted complete graph $K_n$, $\OPT=n-1$ and $LP^*=n/2$:
the all-$1/2$ point is feasible, and summing all edge constraints gives
$(n-1)\sum_vx_v\ge n(n-1)/2$. Hence the worst-case integrality gap is
at least $2-2/n$, approaching 2. This is a family of gap examples, not a
claim that every $n$-vertex graph has this gap.

\subsection{Set cover: bounded frequency}\label{lp:cover}
Let $U$ be a universe, $S_1,\ldots,S_m$ candidate sets with
$\bigcup_jS_j=U$, and costs $c_j\ge0$. One binary variable per set suffices:
\[
  \min\sum_{j=1}^m c_jx_j
  \quad\text{subject to}\quad
  \sum_{j:e\in S_j}x_j\ge1\quad(e\in U),\qquad x_j\in\{0,1\}.
\]
There is no need for an extra variable indicating coverage when every
element must be covered. If each element belongs to at most $k$ sets,
relax and select every set with $x_j^*\ge1/k$. Each covering constraint has
at most $k$ terms, so at least one reaches the threshold. The cost satisfies
\[
  \ALG\le k\sum_jc_jx_j^*=kLP^*\le k\OPT.
\]
This generalizes vertex-cover rounding, where an edge has two endpoint
choices. The hypothesis is bounded \emph{frequency}, not bounded set size.
For example, with one element and two identical cost-1 singleton sets,
$(1/2,1/2)$ is LP-optimal. Each set has size 1, but threshold 1 selects
nothing. This disproves the corresponding threshold proof for bounded set
size.

\subsection{Weighted MAX-SAT: LP rounding}\label{lp:sat}
For the preprocessed formula of Section~\ref{ra:sat}, use $x_j$ for a
variable's truth value and $y_i$ for the reward of clause $C_i$:
\[
  \begin{aligned}
    \max\quad &\sum_iw_i y_i\\
    \text{subject to}\quad &
      y_i\le\sum_{x_j\in C_i}x_j+
             \sum_{\neg x_j\in C_i}(1-x_j)\quad\text{for every }i,\\
    &x_j,y_i\in\{0,1\}.
  \end{aligned}
\]
A false clause forces $y_i=0$; a true clause permits $y_i=1$.
Nonnegative weights make the ILP optimal value exactly the maximum
satisfied weight. Relax all variables to $[0,1]$ and write
$LP^*=\sum_iw_i y_i^*\ge\OPT$.
Independently set $x_j$ true with probability $x_j^*$.

\subsubsection*{One-clause bound: product, AM--GM, then concavity}
For a clause of length $\ell\ge1$, let $a_1,\ldots,a_\ell$ be its
literal-success probabilities: $x_j^*$ for a positive literal and
$1-x_j^*$ for a negative one. Its variables are distinct, and
$\sum_h a_h\ge y_i^*$. Hence
\[
  \begin{aligned}
    \Pr[C_i\text{ satisfied}]
      &=1-\prod_{h=1}^{\ell}(1-a_h)\\
      &\ge1-\left(1-\frac{\sum_h a_h}{\ell}\right)^{\ell}
       \ge1-\left(1-\frac{y_i^*}{\ell}\right)^{\ell}.
  \end{aligned}
\]
The exponent is $\ell$, \emph{not} $1/\ell$.
Let $g_\ell(y)=1-(1-y/\ell)^\ell$. For $\ell=1$ it is linear; for
$\ell\ge2$,
\[
  g_\ell''(y)=-\frac{\ell-1}{\ell}(1-y/\ell)^{\ell-2}\le0.
\]
Since $g_\ell(0)=0$, concavity on $[0,1]$ gives
$g_\ell(y)\ge y g_\ell(1)$. Also
$(1-1/\ell)^\ell\le e^{-1}$. Thus, with
$b_\ell=1-(1-1/\ell)^\ell$,
\[
  \Pr[C_i\text{ satisfied}]
    \ge b_\ell y_i^*\ge(1-1/e)y_i^*,\qquad
  \mathbb E[\ALG]\ge(1-1/e)LP^*\ge(1-1/e)\OPT.
\]
The LP can be solved in polynomial time; sampling and evaluating the result
takes $O(n+L)$ time. The conditional-product calculation in
Section~\ref{ra:conditional} derandomizes the algorithm.

\subsubsection*{Combine fair assignment and LP rounding: $3/4$}
For a clause of length $\ell$, fair assignment satisfies it with probability
$a_\ell=1-2^{-\ell}\ge a_\ell y_i^*$, since $0\le y_i^*\le1$.
LP rounding gives probability at least $b_\ell y_i^*$. The useful comparison
is clause by clause:
\[
  \begin{array}{c|c|c|c}
    \ell&a_\ell&b_\ell&a_\ell+b_\ell\\\hline
    1&1/2&1&3/2\\
    2&3/4&3/4&3/2\\
    \ell\ge3&\ge7/8&\ge1-1/e>5/8&>3/2
  \end{array}
\]
Run both algorithms and return the better assignment. If their values are
$A,B$, then pointwise $\max\{A,B\}\ge(A+B)/2$, so
\[
  \mathbb E[\max\{A,B\}]
    \ge\tfrac12\sum_iw_i(a_{\ell_i}+b_{\ell_i})y_i^*
    \ge\tfrac34LP^*\ge\tfrac34\OPT.
\]
The two outputs need not be independent. Averaging the global constants
$1/2$ and $1-1/e$ would not prove $3/4$; the complementary clause-length
bounds are essential. Alternatively derandomize each algorithm first and
take the better: the same $3/4$ guarantee then holds deterministically.

\subsection{MAX-DICUT: two $1/2$ rounding proofs}\label{lp:dicut}
Let $G=(V,A)$ be a directed graph with nonnegative arc weights and no
self-loops. Put $x_i=1$ when vertex $i$ is on the source side. For each arc,
use a reward variable $z_{ij}$. The ILP and its relaxation are
\[
  \begin{aligned}
    \max\quad&\sum_{(i,j)\in A}w_{ij}z_{ij}\\
    \text{subject to}\quad&z_{ij}\le x_i,\qquad z_{ij}\le1-x_j,\\
    &x_i,z_{ij}\in\{0,1\}\quad
       \text{(replace by }[0,1]\text{ for the LP).}
  \end{aligned}
\]
The upper bounds allow a positive reward only when $x_i=1,x_j=0$.
For any fixed integral cut, setting $z_{ij}=1$ on its outgoing arcs attains
the cut weight, so the ILP is exact. If exact equality
$z_{ij}=x_i(1-x_j)$ is required at every integral feasible point, add
$z_{ij}\ge x_i-x_j$; it is unnecessary for this nonnegative maximization
objective. Below, $x,z$ denote an optimal solution of the displayed
relaxation and $L=LP^*\ge\OPT$.

\subsubsection*{Shifted probabilities: a per-arc comparison}
Independently select vertex $i$ with probability $p_i=1/4+x_i/2$.
For each arc,
\[
  \begin{aligned}
    \Pr[i\in S,j\notin S]
      &=(1/4+x_i/2)(1/4+(1-x_j)/2)\\
      &\ge(1/4+z_{ij}/2)^2\\
      &=z_{ij}/2+(2z_{ij}-1)^2/16\ge z_{ij}/2.
  \end{aligned}
\]
Summing gives $\mathbb E[\ALG]\ge L/2\ge\OPT/2$.
This per-arc inequality holds for any feasible LP point, though optimality
(or another appropriate lower bound on its value) is needed to compare
with $\OPT$.

\subsubsection*{Direct probabilities: a global comparison}
Rounding instead with $p_i=x_i$ also gives a $1/2$ guarantee for an
\emph{optimal} solution of this LP. Let $W=\sum_{(i,j)\in A}w_{ij}$.
If $W=0$, any cut is optimal. Otherwise,
\[
  \Pr[i\in S,j\notin S]=x_i(1-x_j)\ge z_{ij}^2.
\]
Weighted Cauchy--Schwarz gives
\[
  L^2=\left(\sum_{ij}\sqrt{w_{ij}}z_{ij}\sqrt{w_{ij}}\right)^2
    \le\left(\sum_{ij}w_{ij}z_{ij}^2\right)W.
\]
The feasible point $x_i=1/2,z_{ij}=1/2$ has value $W/2$, so optimality
implies $L\ge W/2$. Consequently
\[
  \mathbb E[\ALG]\ge\sum_{ij}w_{ij}z_{ij}^2
    \ge L^2/W\ge L/2\ge\OPT/2.
\]
Direct rounding lacks the shifted scheme's per-arc bound $z_{ij}/2$,
but its global analysis recovers the same factor. Both algorithms are
polynomial time and can be derandomized using conditional expectation.

\subsubsection*{Tuning an affine probability safely}
Let $p_i=a+bx_i$, with $a,b\ge0$ and $a+b\le1$, so all probabilities and
all factors below are nonnegative. Write $c=1-a-b\ge0$. Then
\[
  p_i(1-p_j)\ge(a+bz_{ij})(c+bz_{ij})
    =ac+b(1-b)z_{ij}+b^2z_{ij}^2.
\]
For fixed $b$, $a+c=1-b$ and $ac$ is largest when
$a=c=(1-b)/2$. For this symmetric choice,
\[
  (a+bz)^2=4abz+(a-bz)^2\ge4abz,
  \qquad
  \alpha=4ab=2b(1-b)=\tfrac12-2(b-\tfrac12)^2\le\tfrac12.
\]
Maximize the retained fraction $\alpha$: choose $b=1/2,a=1/4$, obtaining
$1/2$. This optimizes this particular square-based lower bound. It does
not by itself prove that every other analysis of affine rounding is
incapable of a stronger guarantee.

\subsection{Randomized set cover: repeat until feasible}\label{lp:random-cover}
Use the set-cover LP in Section~\ref{lp:cover}, without a frequency bound.
Let $n=|U|\ge1$, $m$ be the number of candidate sets, and
$r=\lceil\ln(2n)\rceil$. (For $U=\varnothing$, return the empty cover.)
Solve the LP once. A \emph{trial} performs $r$ independent rounds; in each
round select each set $S_j$ independently with probability $x_j^*$, then
take the union of all selected sets. Check coverage. Discard a failed trial
and start a fresh independent trial until a cover is found.

\textbf{Success probability.}
For any element $e$, the probability it remains uncovered in one round is
\[
  \prod_{j:e\in S_j}(1-x_j^*)
    \le\exp\left(-\sum_{j:e\in S_j}x_j^*\right)\le e^{-1}.
\]
After $r$ independent rounds its failure probability is at most
$e^{-r}\le1/(2n)$. By the union bound, a trial succeeds with probability
$q\ge1/2$. Element-failure events need not be independent.

\textbf{Cost of a trial and cost of the returned trial.}
A set is included in the union with probability
$1-(1-x_j^*)^r\le r x_j^*$. Thus its nonnegative total cost $C$ satisfies
$\mathbb E[C]\le rLP^*$. However, successful trials can cost more than
typical trials. If $A$ is the success event, the returned trial has the
distribution of one trial conditioned on $A$, and
\[
  \mathbb E[C_{\rm returned}]
    =\mathbb E[C\mid A]
    =\frac{\mathbb E[C\mathbf 1_A]}q
    \le\frac{\mathbb E[C]}q
    \le2rLP^*\le2r\OPT.
\]
Fresh independent trials have a geometric stopping time $T$ with
$\Pr[T>t]=(1-q)^t$ and $\mathbb E[T]=1/q\le2$. Hence the algorithm
terminates almost surely, always returns a checked feasible cover, and
takes expected polynomial time. One trial uses $O(rm)$ random decisions
and $O(n+\sum_j|S_j|)$ time to check coverage after marking selected sets.
The cost guarantee is in expectation; it is not a bound on every returned
cover. The factor $1/q$ cannot be dropped without an additional argument,
because cost and success may be correlated.

\end{document}
